\documentclass[10pt,a4paper]{amsart}
\usepackage[margin=19mm]{geometry}
\usepackage{amsmath,amssymb,mathtools}
\usepackage[scr=rsfs,cal=euler]{mathalfa}
\usepackage{xfrac}
\usepackage[unicode,hidelinks,bookmarks=false]{hyperref}
\newcommand{\ie}{i.e.\ }
\newcommand{\cf}{cf.\ }
\newcommand{\resp}{respectively\ }
\newcommand{\Subsection}{\subsection}
\providecommand{\nobreakdash}{\nobreak\hbox{-}}
\setlength{\emergencystretch}{2em}
\multlinegap0pt
\usepackage{bm}
\usepackage{bbm}
\usepackage{dsfont}
\usepackage{xspace}
\usepackage{cancel}
\usepackage{mathtools}
\usepackage{amsthm}
\makeatletter
\@ifundefined{theorem}{\newtheorem{theorem}{Theorem}}{}
\@ifundefined{lemma}{\newtheorem{lemma}[theorem]{Lemma}}{}
\@ifundefined{prop}{\newtheorem{prop}[theorem]{Proposition}}{}
\makeatother
\newcommand{\holder}{H\"older}
\title[A Note on the Sum-Product Problem and the Convex Sumset Prob]{A Note on the Sum-Product Problem and the Convex Sumset Problem}
\author{Adam Cushman}
\date{Source: arXiv:2512.13849v2 (CC BY 4.0)}
\begin{document}
\maketitle

\noindent\emph{Source and licence.} A Note on the Sum-Product Problem and the Convex Sumset Problem. Version arXiv:2512.13849v2 (CC BY 4.0), distributed under the Creative Commons Attribution 4.0 International licence, as recorded in the arXiv licence metadata for this exact version and shown on the abstract page https://arxiv.org/abs/2512.13849v2. Full rights notice: licenses/arxiv-2512.13849-CC-BY-4.0.txt.\par\smallskip

\noindent\textit{Editorial scope.} Counted unit: the Lemma whose source label is \texttt{lem:SP-szt} in the article (source0.tex lines 849--861, source label \texttt{lem:SP-szt}), delivered together with the article's own complete original proof of it (source0.tex lines 863--1067, one \texttt{proof} environment of 205 lines). No mathematics is written, repaired or re-derived: statement and proof are the article's own text line for line. Required context retained uncounted: lem:solymosi-szt (lines 691--699); prop:rs-result (lines 828--843). Every definition, display and label of those lines is retained unchanged. Omitted from this package and not redistributed: 1--690, 700--827, 844--848, 862--862, 1068--1738. Nothing inside a retained excerpt is omitted or altered: every retained body is delivered exactly as the article's lines are written, so no omission receipt of any kind accompanies this package. Citations: the retained excerpts carry no citation command at all, so no bibliography entry of the article is transcribed and no reference is redistributed. Reference adaptation: none was needed and none was made; every reference inside the delivered unit resolves inside it, so no unresolved reference or citation is produced. Printed slips kept literal: no printed word, symbol, hypothesis or number of the counted statement or of its proof was altered. Layout and class: the article's own class is \texttt{amsart} with packages including amsmath, amsfonts, mathtools, amsthm, amssymb, latexsym, extarrows, xcolor, graphicx, mathrsfs, comment, hyperref, url, pict2e; the wrapper is the lane's pinned \texttt{amsart} editorial wrapper at 10pt on A4, which restates the theorem-like environments the retained text needs and adds the article's own macro definitions that the retained text needs (\texttt{\textbackslash holder}), with extra wrapper packages (bm, bbm, dsfont, xspace, cancel, mathtools). The delivered numbering is the unit's own. Assets: no retained span contains a figure environment, a graphics-inclusion command, a PDF-page inclusion, a table or a TikZ picture; no image file is copied into this package, the package declares no asset dependency and no image file is redistributed with it. Licence: version arXiv:2512.13849v2 of this article is distributed under the Creative Commons Attribution 4.0 International licence, as recorded in the arXiv licence metadata for this exact version and shown on the abstract page https://arxiv.org/abs/2512.13849v2. A full rights notice is delivered at licenses/arxiv-2512.13849-CC-BY-4.0.txt. Characters: every retained excerpt is pure ASCII, so the delivered engine, XeLaTeX with its default Latin Modern text font, reports no missing character.
\begin{lemma}\label{lem:solymosi-szt}    
Let \(A,B \subset \mathbb{R}\) be finite, let \(P = A \times B\), and let \(L\) be either

(a) A finite set of lines whose slopes are finite nonzero real numbers

(b) Finitely many translates of a smooth convex curve

Then the number of incidences between \(P\) and \(L\) is \(\ll \left\lvert P \right\rvert ^{\frac{2}{3} }\left\lvert L \right\rvert ^{\frac{2}{3} } + \left\lvert L \right\rvert     \)
\end{lemma}

\begin{prop}\label{prop:rs-result}
Let \(A \subset \mathbb{R}^{+} \). Let \(\tau \in \mathbb{R} \) be so that, defining
\[
    S = \left\{ \lambda \in \frac{A}{A}  : r_{\frac{A}{A} } (\lambda) \in [\tau,2\tau) \right\} 
,\]
we have
\[
    E^{\times }(A) \approx \tau^{2} \left\lvert S  \right\rvert 
.\]

There exists \(S ' \subset S \) with \(\left\lvert S ' \right\rvert \geq \frac{1}{64} \cdot \left\lvert S \right\rvert \) such that,
for all \(\lambda \in S'\),
\[
    \left\lvert A A_{\lambda}  \right\rvert \gtrsim \frac{\left\lvert A \right\rvert ^{18}}{\left\lvert S  \right\rvert ^{\frac{1}{2} } \left\lvert AA \right\rvert ^{4} \left\lvert A + A \right\rvert ^{8}}
.\]
\end{prop}

\begin{lemma}\label{lem:SP-szt}
Let \(A \subset \mathbb{R} ^{+}\). There is \(A_0 \subset A\) with \(\left\lvert A_0 \right\rvert > \frac{1}{2} \cdot \left\lvert A \right\rvert \) for which

(1) For all \(B \subset \mathbb{R} \), 
\[
    E_3(A_0,B) \lesssim \frac{\left\lvert B \right\rvert ^{2} \left\lvert AA \right\rvert ^{\frac{35}{2} } \left\lvert A + A \right\rvert ^{24}}{\left\lvert A \right\rvert ^{54}} 
\]

(2) For all \(B \subset \mathbb{R} \), and \(s \in (1,3)\),
\[
    E_{s} (A_0,B) \lesssim \left\lvert B \right\rvert ^{\frac{1 + s}{2} } \left\lvert A \right\rvert ^{\frac{1}{2} \cdot \left( 57 - 55s \right) } \left\lvert AA \right\rvert ^{\frac{35}{4} (s-1)} \left\lvert A+A \right\rvert ^{12 \left( s-1 \right) }
.\]
\end{lemma}

\begin{proof}
We first prove (1), from which (2) immediately follows by interpolation for \(E_{s} (A,B)\)
between \(E_3(A,B)\) and \(E_{1} (A,B) = \left\lvert A \right\rvert \left\lvert B \right\rvert \).

For the sake of clarity, let \(\Pi = AA\).

Suppose \(X \subset A\) with \(\left\lvert X \right\rvert  > \frac{1}{2} \cdot \left\lvert A \right\rvert \).
We apply Proposition \ref{prop:rs-result} to \(X\) instead of \(A\).
We show that, given \(X\) as above, Proposition \ref{prop:rs-result} gives rise to a
nonempty set \(X'\), a subset of some dilate of \(X\), for which
\(r_{\Pi / \Pi} (\lambda)\) is large for \(\lambda \in X'\). 

From this, we obtain
a bound on \(E_3(X',B)\). We then iterate this
general result, beginning with \(A\), obtaining a sequence of sets \(\left\{ A_{i} ' \right\} \),
and a bound on each of \(E_3(A_{i} ',B)\). The set \(\cup _{i} A_{i}' \) will be our \(A_0\).

A dyadic partitioning gives \(\beta \in \mathbb{R} \) so that, defining
\[
    S  : = \left\{ \lambda \in \frac{X}{X} : r_{\frac{X}{X} } (\lambda) \in [\beta,2\beta) \right\} 
\]
we have
\[
    E^{\times }(X) \approx \beta ^{2}\left\lvert S  \right\rvert 
.\]
Proposition \ref{prop:rs-result} applied to \(X\) yields \(S ' \subset S \) with \(\left\lvert S ' \right\rvert \geq \frac{1}{64} \cdot \left\lvert S  \right\rvert \) for which
for any \(\lambda \in S'\),
\begin{equation} \label{eq:xxlambda-bound}
    \left\lvert X X_{\lambda}  \right\rvert \gtrsim \frac{\left\lvert X \right\rvert ^{18}}{\left\lvert S \right\rvert ^{\frac{1}{2} } \left\lvert XX \right\rvert ^{4} \left\lvert X + X \right\rvert ^{8}}.
\end{equation}  
Since \(X \subset A\), \(\left\lvert A A_{\lambda}  \right\rvert \geq \left\lvert X X_{\lambda}  \right\rvert \). Additionally, \(\left\lvert X \right\rvert \gg \left\lvert A \right\rvert \)
and \(\left\lvert X + X \right\rvert \leq \left\lvert A+A \right\rvert \), \(\left\lvert XX \right\rvert \leq \left\lvert \Pi \right\rvert \).
Substituting all of these into (\ref{eq:xxlambda-bound}) we have that for any \(\lambda \in S'\),
\[
    \left\lvert A A_{\lambda}  \right\rvert \gtrsim  \frac{\left\lvert A \right\rvert ^{18}}{\left\lvert S \right\rvert ^{\frac{1}{2} } \left\lvert \Pi \right\rvert ^{4} \left\lvert A + A \right\rvert ^{8}} 
.\]

Observe the following truism. For \(\lambda \in S',~ a \in A,~ a_{\lambda} \in A_{\lambda} \). We have
\[
    \lambda = \frac{a \left( \lambda a_{\lambda}  \right) }{a a_{\lambda} } \in \frac{\Pi}{\Pi} 
.\]
As such, each distinct \(d \in A A_{\lambda} \) gives rise to the representation \((\lambda d, d) \in \Pi ^{2}\), and hence, \(\forall \lambda \in S '\),
\[
    r_{ \Pi / \Pi  } (\lambda) \geq \left\lvert A A_{\lambda}  \right\rvert \gtrsim \frac{\left\lvert A \right\rvert ^{18}}{\left\lvert S \right\rvert ^{\frac{1}{2} } \left\lvert \Pi \right\rvert ^{4} \left\lvert A + A \right\rvert ^{8}} 
.\]

We pass from elements in \(S'\) to elements in some dilate of \(X\) by pigeonholing.
Using the definition of \(S'\) and the definition of \(r_{X / X} \) respectively, we have
\[
    \beta \left\lvert S' \right\rvert \leq  \sum _{\lambda \in S'} r_{ X / X } (\lambda) = \sum _{\lambda \in S'} \sum _{x_0 \in X} \# \left\{ x \in X : \frac{x}{x_0} = \lambda \right\}   = \sum _{x_0 \in X} \# \left\{ x \in X : \frac{x}{x_0} \in S ' \right\}  
,\]
and hence there is \(x_0 \in X\) such that
\[
    \# \left\{ x \in X : \frac{x}{x_0} \in S'  \right\} \geq \frac{\beta \left\lvert S  \right\rvert }{64 \left\lvert X \right\rvert } 
.\]
Letting
\[
    X' = \left( \frac{X}{x_0}  \right) \cap S'
,\]
we see that \(\left\lvert X' \right\rvert \gg \frac{\beta \left\lvert S  \right\rvert }{\left\lvert X \right\rvert } \), and
for any \(x \in X'\)
\begin{equation} \label{eq:xprime-rep}
    r_{ \Pi / \Pi } \left( x  \right)  \gtrsim \frac{\left\lvert A \right\rvert ^{18}}{\left\lvert S \right\rvert ^{\frac{1}{2} } \left\lvert \Pi \right\rvert ^{4} \left\lvert A + A \right\rvert ^{8}}.
\end{equation}  

We have
\[
    \left\lvert X' \right\rvert \gg \frac{\beta \left\lvert S \right\rvert }{\left\lvert X \right\rvert }  \implies \beta \left\lvert S \right\rvert \ll \left\lvert X \right\rvert \left\lvert X' \right\rvert \leq \left\lvert A \right\rvert \left\lvert X' \right\rvert 
,\]
and therefore, substituting into the RHS of (\ref{eq:xprime-rep}),
\[
    \frac{\left\lvert A \right\rvert ^{18}}{\left\lvert S \right\rvert ^{\frac{1}{2} } \left\lvert \Pi \right\rvert ^{4} \left\lvert A + A \right\rvert ^{8}} \gg \frac{\left\lvert A \right\rvert ^{17} \beta \left\lvert S \right\rvert ^{\frac{1}{2} }}{\left\lvert X' \right\rvert \left\lvert \Pi \right\rvert ^{4} \left\lvert A + A \right\rvert ^{8}}
.\]
By the definition of \(\beta , S\), Cauchy-Schwarz, and \(\left\lvert X \right\rvert \gg \left\lvert A \right\rvert \), \(\left\lvert X X \right\rvert \leq \left\lvert \Pi \right\rvert \) respectively, we have
\[
    \beta ^{2} \left\lvert S \right\rvert \approx E^{\times }(X) \geq \frac{\left\lvert X \right\rvert ^{4}}{\left\lvert XX \right\rvert }  \gg \frac{\left\lvert A \right\rvert ^{4}}{\left\lvert \Pi \right\rvert } 
,\]
and hence
\[
    \frac{\left\lvert A \right\rvert ^{17} \beta \left\lvert S \right\rvert ^{\frac{1}{2} }}{\left\lvert X' \right\rvert \left\lvert \Pi \right\rvert ^{4} \left\lvert A + A \right\rvert ^{8}} \gg \frac{\left\lvert A \right\rvert ^{17} \left( \frac{\left\lvert A \right\rvert ^{4}}{\left\lvert \Pi \right\rvert }  \right) ^{\frac{1}{2} }}{\left\lvert X' \right\rvert \left\lvert \Pi \right\rvert ^{4} \left\lvert A + A \right\rvert ^{8}} = \frac{1}{\left\lvert X' \right\rvert } \cdot \frac{\left\lvert A \right\rvert ^{19}}{\left\lvert \Pi \right\rvert ^{\frac{9}{2} }\left\lvert A+A \right\rvert ^{8}}  
.\]

Substituting into (\ref{eq:xprime-rep}) gives that for any \(x \in X'\) we have
\begin{equation} \label{eq:xprime-rep-final}
    r_{ \Pi / \Pi } \left( x  \right)  \gtrsim \frac{1}{\left\lvert X' \right\rvert } \cdot \frac{\left\lvert A \right\rvert ^{19}}{\left\lvert \Pi \right\rvert ^{\frac{9}{2} } \left\lvert A + A \right\rvert ^{8}} .
\end{equation}

We see that \(\left\lvert X' \right\rvert \geq 1\), or else
\[
    \beta \left\lvert S \right\rvert \ll \left\lvert X \right\rvert \leq \left\lvert A \right\rvert 
,\]
and hence, dividing from
\[
    \beta^{2} \left\lvert S \right\rvert \approx E^{\times }(X) \geq \frac{\left\lvert X \right\rvert ^{4}}{\left\lvert XX \right\rvert } \gg \frac{\left\lvert A \right\rvert ^{4}}{\left\lvert \Pi \right\rvert } 
,\]
we see that
\[
    \beta \gtrsim \frac{\left\lvert A \right\rvert ^{3} }{\left\lvert \Pi \right\rvert } 
.\]
Using the trivial bound \(\beta \leq \left\lvert A \right\rvert \) gives
\(\left\lvert \Pi \right\rvert \gtrsim \left\lvert A \right\rvert ^{2}\), and hence Lemma \ref{lem:SP-szt} part 1 holds trivially
upon taking \(A_0 = A\).

We proceed by obtaining an energy bound for \(X'\).
Fix an arbitrary \(B \subset \mathbb{R} \) finite.

By a dyadic partitioning, there is \(k \in \mathbb{N} \) and
\[
    D^{(k)} = \left\{ d \in X'  - B  : \delta_{X',B} (d) \in [k,2k) \right\}
\]
such that
\[
    E_{3} (X',B)\approx k ^{3} \left\lvert D^{(k)} \right\rvert 
.\]

By the definition of \(D^{(k)}\),
\[
    \sum _{d \in D^{(k)}} \delta_{X',B} (d) \geq k \left\lvert D^{(k)} \right\rvert 
.\]
We have
\[
    \sum _{d \in D^{(k)}} \delta_{X',B} (d) = \sum _{x \in X'} \sigma_{B,D^{(k)}} (x)
,\]
and hence, by (\ref{eq:xprime-rep-final}), we have
\begin{equation} \label{eq:lem-lower-bound}
    \sum _{x \in X'} \sigma_{B,D^{(k)}} (x) r_{ \Pi / \Pi} (x) \gtrsim k \left\lvert D^{(k)} \right\rvert \cdot \frac{1}{\left\lvert X' \right\rvert } \cdot \frac{\left\lvert A \right\rvert ^{19}}{\left\lvert \Pi \right\rvert ^{\frac{9}{2} } \left\lvert A + A \right\rvert ^{8}}  .
\end{equation}  

Trivially we have
\[
    \sum _{x \in X'} \sigma_{B,D^{(k)}} (x) r_{ \Pi / \Pi} (x) \leq \sum _{x} \sigma_{B,D^{(k)}} (x) r_{ \Pi / \Pi} (x)
.\]
See that
\[
    \sum _{x} \sigma_{B,D^{(k)}} (x) r_{ \Pi / \Pi} (x)
\]
counts solutions to
\begin{equation} \label{eq:point-line-system}
    \frac{1}{\pi_2} \cdot \pi_1 - d = b,~ \pi_{i} \in \Pi,~ b \in B,~ d \in D^{(k)}.
\end{equation}
Solutions to (\ref{eq:point-line-system}) are precisely the incidences of the point set
\(\Pi \times B\) with the system of lines
\[
    \ell(x) = \frac{x}{\pi} - d ,~ \pi \in \Pi ,~ d \in D^{(k)}
.\]
Since \(\Pi \subset \mathbb{R} _{>0} \), the slopes of
these lines are finite nonzero real numbers, so applying Lemma \ref{lem:solymosi-szt} gives
\[
    \sum _{x} \sigma_{B,D^{(k)}} (x) r_{ \Pi / \Pi} (x) \ll \left( \left\lvert \Pi \right\rvert ^{2} \left\lvert B \right\rvert \left\lvert D^{(k)} \right\rvert  \right) ^{\frac{2}{3} } + \left\lvert \Pi \right\rvert \left\lvert D^{(k)} \right\rvert 
.\]
Using the trivial bound \(\left\lvert D^{(k)} \right\rvert \leq  \left\lvert A \right\rvert \left\lvert B \right\rvert \) we have
\[
    \left\lvert \Pi \right\rvert \left\lvert D^{(k)} \right\rvert \ll \left( \left\lvert \Pi \right\rvert ^{2} \left\lvert B \right\rvert \left\lvert D^{(k)} \right\rvert  \right) ^{\frac{2}{3} }
,\]
so combining with (\ref{eq:lem-lower-bound}) gives
\[
    k \left\lvert D^{(k)} \right\rvert \cdot \frac{1}{\left\lvert X' \right\rvert } \cdot \frac{\left\lvert A \right\rvert ^{19}}{\left\lvert \Pi \right\rvert ^{\frac{9}{2} } \left\lvert A + A \right\rvert ^{8}} \lesssim \sum _{x} \sigma_{B,D^{(k)}} (x) r_{ \Pi / \Pi} (x) \ll \left( \left\lvert \Pi \right\rvert ^{2} \left\lvert B \right\rvert \left\lvert D^{(k)} \right\rvert  \right) ^{\frac{2}{3} }
,\]
or equivalently
\[
    k^{3}  \left\lvert D^{(k)} \right\rvert \lesssim \left\lvert X' \right\rvert ^{3} \cdot \frac{\left\lvert B \right\rvert ^{2} \left\lvert \Pi \right\rvert ^{\frac{35}{2} }\left\lvert A+A \right\rvert ^{24}}{\left\lvert A \right\rvert ^{57}} 
.\]
Recall that \(k ^{3} \left\lvert D^{(k)} \right\rvert \approx E_3(X',B)\).

Having found the desired bound on \(E_3(X',B)\) we proceed with iteration.
Let \(A_1 = A\). 
If \(\left\lvert A_{j} \right\rvert  > \frac{1}{2} \cdot \left\lvert A \right\rvert  \), obtain \(A_{j+1} \) from \(A_{j} \) by applying the above argument
with \(X = A_{j} \), yielding some \(a_{j}  \in A\) and a subset \(A_{j} ' \subset A_{j} \) for which we have
\[
    \forall a \in A_{j} ' ,~ r_{ \Pi / \Pi } \left( \frac{a}{a_{j} }   \right)  \gtrsim \frac{1}{\left\lvert A_{j} ' \right\rvert } \cdot \frac{\left\lvert A \right\rvert ^{19}}{\left\lvert \Pi \right\rvert ^{\frac{9}{2} } \left\lvert A + A \right\rvert ^{8}} 
,\]
and hence for any \(B \subset \mathbb{R} \),
\[
    E_3(A_{j} ',B) = E_3 \left( \frac{A_{j} '}{a_{j} } , \frac{B}{a_{j} }    \right) \lesssim \left\lvert A_{j} ' \right\rvert ^{3} \cdot \frac{\left\lvert B \right\rvert ^{2} \left\lvert \Pi \right\rvert ^{\frac{35}{2} }\left\lvert A+A \right\rvert ^{24}}{\left\lvert A \right\rvert ^{57}} 
.\]
Let \(A_{j+1} = A_{j} \setminus A_{j} ' \). Once \(\left\lvert A_{n}  \right\rvert \leq \frac{1}{2} \cdot \left\lvert A \right\rvert     \), which will occur in a finite number of steps since \(\left\lvert A_{j} '  \right\rvert \geq 1\),
let
\[
    A_0 = \bigsqcup  _{j=1} ^{n-1}A_{j} ' = A \setminus A_{n} 
.\]
We have \(\left\lvert A_0 \right\rvert > \frac{1}{2} \cdot \left\lvert A \right\rvert \) and, by Minkowski's inequality, for any \(B \subset \mathbb{R} \),
\[
    E_3(A_0,B)^{\frac{1}{3} } = \left\lVert \delta_{A_0,B}  \right\rVert _{3} = \left\lVert \sum _{j =1} ^{n-1} \delta_{A_{j}' ,B}  \right\rVert _{3} \leq \sum _{j=1} ^{n-1} \left\lVert \delta_{A_{j}' ,B}  \right\rVert _{3} \lesssim \frac{\left\lvert B \right\rvert ^{\frac{2}{3} } \left\lvert \Pi \right\rvert ^{\frac{35}{6}  } \left\lvert A+ A \right\rvert ^{8}}{\left\lvert A \right\rvert ^{19}} \sum _{j=1} ^{n-1} \left\lvert A_{j} ' \right\rvert 
.\]
Seeing that \( \sum _{j=1} ^{n-1}\left\lvert A_{j} ' \right\rvert \leq \left\lvert A \right\rvert   \) gives
\[
    E_3(A_0,B)\lesssim \frac{\left\lvert B \right\rvert^{2}  \left\lvert \Pi \right\rvert ^{\frac{35}{2} } \left\lvert A+A \right\rvert ^{24}}{\left\lvert A \right\rvert ^{54}} 
\]
as desired. Part (1) is complete.

For part (2), fix \(s \in (1,3)\) and use \holder's inequality to get
\[
    \sum _{x} \delta_{A_0,B} (x)^{s} = \sum _{x} \delta_{A_0,B} (x) ^{3 \cdot \frac{s-1}{2}  } \delta_{A_0,B} (x) ^{\frac{3- s }{2} } \leq \left( \sum _{x} \delta_{A_0,B} (x)^{3}  \right) ^{\frac{s-1}{2} } \left( \sum _{x} \delta_{A_0,B} (x) \right) ^{\frac{3-s}{2} }
.\]
Seeing that, by part (1) and the trivial result \( \sum _{x} \delta_{A_0,B} (x) = \left\lvert A_0 \right\rvert \left\lvert B \right\rvert \),
we have
\[
    \left( \sum _{x} \delta_{A_0,B} (x)^{3}  \right) ^{\frac{s-1}{2} } \left( \sum _{x} \delta_{A_0,B} (x) \right) ^{\frac{3-s}{2} } \lesssim \left( \frac{\left\lvert B \right\rvert^{2}  \left\lvert \Pi \right\rvert ^{\frac{35}{2} } \left\lvert A+A \right\rvert ^{24}}{\left\lvert A \right\rvert ^{54}}  \right) ^{\frac{s-1}{2} }\left( \left\lvert A_0 \right\rvert \left\lvert B \right\rvert  \right) ^{\frac{3-s}{2} }
,\]
and so
\[
    E_{s} (A_0,B) \lesssim \left\lvert B \right\rvert ^{\frac{1 + s}{2} } \left\lvert A \right\rvert ^{\frac{1}{2} \cdot \left( 57 - 55s \right) } \left\lvert \Pi \right\rvert ^{\frac{35}{4} (s-1)} \left\lvert A+A \right\rvert ^{12 \left( s-1 \right) }
.\]

\end{proof}

\end{document}
